\chapter{АНАЛИЗ ПРЕДМЕТНОЙ ОБЛАСТИ}

% -------------------------------------------------------
\section{Когнитивные предубеждения: теоретические основания}
% -------------------------------------------------------

\subsection{Программа эвристик и предубеждений Канемана и Тверски}

Систематическое изучение когнитивных предубеждений как научной проблемы восходит к серии работ Амоса Тверски и Даниэля Канемана, опубликованных в 1970--1980-х годах~\cite{tversky1981framing, tversky1983extensional}. В основе их исследовательской программы лежало наблюдение о том, что люди при принятии решений в условиях неопределённости систематически нарушают нормы рационального поведения, постулируемые теорией ожидаемой полезности и математической теорией вероятностей. Вместо точного байесовского обновления убеждений человек использует упрощённые ментальные операции --- эвристики, --- которые, будучи когнитивно экономными, порождают предсказуемые и воспроизводимые ошибки суждения.

Канеман и Тверски описали три основные эвристики, лежащие в основе широкого класса предубеждений~\cite{kahneman2011thinking}. \textit{Эвристика репрезентативности} состоит в оценке вероятности события по степени его соответствия типичному прототипу или нарративному шаблону без учёта базовых вероятностей. \textit{Эвристика доступности} связана с оценкой вероятности через лёгкость, с которой соответствующие примеры извлекаются из памяти. \textit{Эвристика якорения и корректировки} предполагает оценку числовых величин через корректировку от некоторого начального значения, причём корректировка систематически оказывается недостаточной.

Для настоящего исследования наиболее релевантна эвристика репрезентативности, поскольку именно она лежит в основе ошибки конъюнкции~--- центральной парадигмы первого эксперимента. Согласно этой эвристике, конъюнктивный вариант ответа кажется более вероятным, если он нарративно согласуется с описанием персонажа, даже если формальные законы теории вероятностей однозначно исключают такой вывод.

Теоретическим обобщением эвристической программы стала двухсистемная модель, разработанная Канеманом~\cite{kahneman2011thinking}: \textit{Система 1} работает быстро, автоматически, ассоциативно и не требует осознанных усилий; \textit{Система 2} --- медленно, последовательно, по правилам и с затратами когнитивных ресурсов. Большинство когнитивных предубеждений возникает именно тогда, когда Система 1 генерирует ответ, который Система 2 принимает без достаточной проверки. Эта модель имеет прямые следствия для интерпретации поведения языковых моделей: авторегрессионная генерация токенов по своей природе ближе к Системе 1, тогда как явное пошаговое рассуждение (chain-of-thought) представляет собой попытку эмулировать Систему 2~\cite{wei2022chain}.

\subsection{Типология когнитивных предубеждений: вероятностные, логические, индуктивные}

Для целей настоящего исследования когнитивные предубеждения классифицируются по типу нарушаемой нормы рассуждения. Эта классификация непосредственно определяет выбор трёх экспериментальных парадигм и позволяет рассматривать их как единый спектр, а не как изолированные явления.

\textit{Вероятностные предубеждения} возникают при нарушении аксиом теории вероятностей. Наиболее известный пример --- ошибка конъюнкции~\cite{tversky1983extensional}: оценка $P(A \cap B) > P(A)$, что математически невозможно. К этой же категории относятся пренебрежение базовыми вероятностями (base-rate neglect), ошибка игрока и эффект фрейминга~\cite{tversky1981framing}. Общим механизмом является подмена вероятностного суждения оценкой нарративной правдоподобности.

\textit{Логические предубеждения} проявляются при нарушении норм дедуктивного вывода. Задача выбора карточек Уэйсона~\cite{wason1966reasoning} является классическим инструментом их диагностики: правильный ответ требует применения правила modus tollens ($P \rightarrow Q$, $\lnot Q \Rightarrow \lnot P$), тогда как большинство испытуемых выбирают подтверждающие карточки вместо фальсифицирующих. Предубеждение подтверждения~\cite{nickerson1998confirmation} --- тенденция искать и интерпретировать информацию в пользу имеющейся гипотезы --- является общим механизмом как для логических, так и для индуктивных предубеждений.

\textit{Индуктивные предубеждения} возникают при нарушении нормативных стандартов формирования и пересмотра гипотез. Задача 2-4-6 Уэйсона~\cite{wason1960failure} демонстрирует, что люди систематически тестируют тройки, согласующиеся с текущей гипотезой, вместо того чтобы искать опровергающие примеры. Байесовская норма предписывает максимизацию информационной ценности теста, что требует проверки именно тех случаев, которые могут фальсифицировать гипотезу~\cite{tenenbaum1999bayesian}.

Три парадигмы настоящего исследования образуют прогрессирующий спектр сложности рассуждения: от одношагового вероятностного суждения к дедуктивному анализу условных правил и далее к многоходовому итеративному открытию правил. Это позволяет исследовать, является ли предубеждение LLM единым феноменом или варьируется в зависимости от типа требуемого рассуждения.

\subsection{Нормативные стандарты рассуждения и отклонения от них}

Под нормативным стандартом в настоящей работе понимается критерий, относительно которого ответ модели однозначно классифицируется как правильный или ошибочный. Для каждой из трёх парадигм такой стандарт задаётся формально.

В задаче ошибки конъюнкции нормативный ответ определяется аксиомой теории вероятностей: $P(A \cap B) \leq P(A)$ для любых событий $A$ и $B$. Выбор конъюнктивного варианта однозначно является ошибкой вне зависимости от нарративного контекста.

В задаче выбора карточек нормативный ответ определяется правилом материальной условности: для проверки истинности правила «если $P$, то $Q$» необходимо и достаточно перевернуть карточки $P$ и $\lnot Q$~\cite{wason1966reasoning}. Карточки $\lnot P$ и $Q$ не могут опровергнуть правило и поэтому не должны выбираться.

В задаче 2-4-6 нормативный стандарт является более сложным: критерием успеха служит финальная гипотеза модели, а не конкретная последовательность тестов. Оценка производится через метрику пересечения по объединению (IOU) над случайной выборкой троек, что позволяет сравнивать гипотезы, записанные в различных синтаксических формах, по их экстенсиональной эквивалентности~\cite{banatt2024wilt}.

% -------------------------------------------------------
\section{Анализ проблемы: когнитивные предубеждения в LLM}
% -------------------------------------------------------

\subsection{Природа предубеждений в языковых моделях: memorization vs. emergent bias}

Принципиальный методологический вопрос при тестировании LLM на когнитивных задачах состоит в интерпретации наблюдаемых ошибок. Существуют два конкурирующих объяснения того, почему языковая модель может демонстрировать ответ, характерный для когнитивного предубеждения.

Согласно гипотезе \textit{memorization}, модель воспроизводит правильный или ошибочный ответ, потому что видела данную конкретную задачу или её близкие аналоги в обучающем корпусе. Задача Линды, задача Уэйсона и задача 2-4-6 входят в число наиболее цитируемых психологических экспериментов и с высокой вероятностью представлены в корпусах, на которых обучены современные LLM. Если модель демонстрирует правильный ответ на каноническую формулировку, это не обязательно свидетельствует о подлинном вероятностном рассуждении --- она могла просто запомнить ответ~\cite{zhou2023benchmark}. Симметрично, воспроизведение человекоподобной ошибки на знакомой задаче может отражать memorization паттернов ошибок из психологических текстов, а не возникновение bias в процессе рассуждения.

Согласно гипотезе \textit{emergent bias}, предубеждения возникают как следствие самой природы обучения на человеческих текстах: модель усваивает не только факты, но и типичные паттерны человеческого рассуждения, включая систематические отклонения от нормативных стандартов. В этом случае предубеждения должны воспроизводиться на новых, контаминационно чистых стимулах, структурно аналогичных классическим задачам.

Binz и Schulz~\cite{binz2023using} продемонстрировали свидетельства в пользу гипотезы memorization: GPT-3 решал каноническую версию задачи Уэйсона правильно, но немедленно давал ошибочные ответы при изменении порядка карточек. Dasgupta и соавторы~\cite{dasgupta2024language}, напротив, показали, что на специально сконструированных деконтаминированных стимулах крупные модели воспроизводят человекоподобные контент-эффекты, что указывает на emergent bias. Настоящее исследование придерживается осторожной позиции: контаминация контролируется через генерацию оригинальных датасетов, однако разграничение memorization и emergent bias не является первичной целью --- целью является измерение итогового поведения модели вне зависимости от механизма его порождения.

\subsection{Проблема контаминации обучающих данных}

Контаминация обучающих данных (training data contamination) --- одна из центральных методологических проблем при оценке когнитивных способностей LLM~\cite{zhou2023benchmark}. Её специфика в контексте когнитивных задач состоит в следующем: в отличие от математических бенчмарков, где контаминация ведёт к завышению точности, в когнитивных задачах контаминация может давать как завышение (если модель запоминает правильный нормативный ответ), так и занижение (если модель запоминает человекоподобный ошибочный паттерн из психологических учебников).

Предшествующие исследования предлагали различные стратегии борьбы с контаминацией. Suri и соавторы~\cite{suri2024large} разработали оригинальный промпт с новым сценарием (посетительница кафе), исключающим прямое совпадение с текстом задачи Линды. Jiang и соавторы~\cite{jiang2024peek} предложили фреймворк динамической генерации синтетических задач на основе корпуса SCT~\cite{mostafazadeh2016corpus}, что исключает возможность присутствия конкретных комбинаций в обучающих данных. Chadimova и соавторы~\cite{chadimova2024meaningless} показали, что замена смысловых слов в промпте хеш-значениями улучшает нормативную точность, что интерпретируется как свидетельство токен-bias. Arkoudas~\cite{arkoudas2023gpt4} продемонстрировал, что GPT-4 не справляется с поверхностно модифицированными версиями логических задач.

В настоящем исследовании контаминация контролируется на уровне каждого эксперимента. Для задачи Линды используется оригинальный датасет из 110 виньеток с системой демографических подстановок. Для задачи Уэйсона~--- четыре стимульных датасета, из которых три включают нека нонические и оригинальные стимулы. Для задачи 2-4-6~--- независимый датасет из 50 правил с верификацией на отсутствие функциональных дубликатов.

\subsection{Ограничения существующих подходов к тестированию}

Анализ литературы позволяет выделить несколько системных ограничений, характерных для существующих исследований когнитивных предубеждений в LLM.

\textit{Узость охвата моделей.} Большинство работ тестируют одну-две модели, как правило из семейства GPT~\cite{binz2023using, suri2024large}, что не позволяет делать выводы о том, являются ли наблюдаемые паттерны специфическими для конкретной архитектуры или универсальными.

\textit{Бинарность метрик.} Преобладающим подходом является подсчёт доли правильных/ошибочных ответов~\cite{macmillan2024irrationality, jiang2024peek}. Эта метрика не различает модель, дающую правильный ответ с уверенностью 51\%, и модель, дающую его с уверенностью 99\%. Между тем уверенность модели в ошибочном ответе является самостоятельной характеристикой силы предубеждения.

\textit{Изоляция парадигм.} Каждая работа, как правило, фокусируется на одной задаче и не рассматривает предубеждения как феномен, существующий на спектре~\cite{yax2024studying}. Это не позволяет исследовать, связаны ли предубеждения в разных парадигмах на уровне отдельных моделей.

\textit{Отсутствие анализа стабильности.} Большинство работ сообщают результаты по одной версии промпта и одному порядку предъявления стимулов, не измеряя вариабельность ответов модели как самостоятельную переменную~\cite{macmillan2024irrationality}.

Настоящее исследование адресует каждое из перечисленных ограничений: тестируются девять моделей, вводится непрерывная метрика силы предубеждения, исследуются три парадигмы в едином дизайне, а стабильность ответов измеряется через систематическую перестановку стимулов.

% -------------------------------------------------------
\section{Обзор смежных исследований}
% -------------------------------------------------------

\subsection{Ошибка конъюнкции и эвристика репрезентативности в LLM}

Задача Линды была одной из первых классических психологических парадигм, применённых к тестированию языковых моделей. Binz и Schulz~\cite{binz2023using} включили её в свой систематический бенчмарк GPT-3 и обнаружили, что модель воспроизводит ошибку конъюнкции в условиях, аналогичных человеческому эксперименту. Авторы интерпретировали этот результат как свидетельство того, что GPT-3 использует эвристику репрезентативности.

Suri и соавторы~\cite{suri2024large} расширили эту линию исследований на GPT-3.5 и показали, что модель демонстрирует ошибку конъюнкции с частотой, сопоставимой с человеческой ($\approx$80\%), что указывает на воспроизведение нормативно некорректного паттерна. Wang и соавторы~\cite{wang2024linda} сосредоточились на манипуляции репрезентативностью: они показали, что современные модели (GPT-4, Claude) реже допускают ошибку конъюнкции, чем более ранние, однако сохраняют чувствительность к нарративной согласованности конъюнктивного элемента.

Jiang и соавторы~\cite{jiang2024peek} ввели понятие \textit{токен-bias}: предрасположенности модели выбирать ответ с более высокой априорной вероятностью в языковом корпусе вне зависимости от логической структуры задачи. Для его измерения авторы использовали тест МакНемара на парах задач с коррелированным и некоррелированным конъюнктивным элементом --- подход, непосредственно воспринятый в настоящем исследовании.

Ryu и соавторы~\cite{ryu2024representativeness} и Saeedi и Goodarzi~\cite{saeedi2024gpt} подтвердили устойчивость ошибки конъюнкции в различных моделях и формулировках, в том числе при смене языка промпта. Chadimova и соавторы~\cite{chadimova2024meaningless} предложили интерпретацию через механизм токен-предубеждения: замена нарративно нагруженных слов на нейтральные хеш-значения существенно снижает частоту ошибки.

\subsection{Задача выбора карточек Уэйсона в LLM}

Задача Уэйсона является наиболее изученной логической парадигмой в исследованиях LLM. Binz и Schulz~\cite{binz2023using} показали, что GPT-3 решает каноническую абстрактную версию правильно, однако это, по всей видимости, является следствием memorization: при перестановке карточек точность падает. Arkoudas~\cite{arkoudas2023gpt4} продемонстрировал, что GPT-4 не справляется со структурно эквивалентной задачей при поверхностной модификации стимулов.

Dasgupta и соавторы~\cite{dasgupta2024language} провели наиболее методологически строгое исследование: они использовали специально сконструированные деконтаминированные стимулы и обнаружили, что крупные модели воспроизводят контент-эффекты, аналогичные человеческим~--- социальные контракты решаются точнее, чем абстрактные правила, что согласуется с теорией прагматических схем рассуждения~\cite{cheng1985pragmatic}.

Macmillan-Scott и Musolesi~\cite{macmillan2024irrationality} провели масштабный бенчмарк нескольких моделей и выявили существенную вариабельность: одни модели демонстрируют точность, близкую к потолку, другие --- систематические ошибки предубеждения соответствия (matching bias). Авторы также зафиксировали высокую нестабильность ответов при повторных запросах с идентичными промптами, что поставило вопрос о надёжности точечных оценок.

Seals и Shalin~\cite{seals2024evaluating} исследовали дедуктивную компетентность LLM в расширенном наборе логических задач и пришли к выводу, что современные модели демонстрируют систематическую зависимость от поверхностных признаков стимула, а не от логической структуры. Это согласуется с предубеждением соответствия, описанным Эвансом~\cite{evans1998matching}.

\subsection{Индуктивное открытие правил и задача 2-4-6 в LLM}

Исследования индуктивного рассуждения LLM значительно менее развиты по сравнению с вероятностными и логическими задачами. Banatt и соавторы~\cite{banatt2024wilt} разработали бенчмарк WILT --- первый систематический инструмент оценки LLM в многоходовой задаче открытия правил. Их результаты показывают, что более крупные модели демонстрируют более высокий успех в открытии правил, однако все модели проявляют признаки предубеждения подтверждения: доля подтверждающих тестов превышает нормативную.

Bazigaran и Sohn~\cite{bazigaran2025concept} исследовали концептуальное обобщение в LLM на материале «числовой игры» Тененбаума~\cite{tenenbaum1999bayesian} и обнаружили паттерны, частично соответствующие байесовскому обновлению убеждений. Coda-Forno и соавторы~\cite{coda2024cogbench} включили задачи индуктивного рассуждения в комплексный когнитивный бенчмарк CogBench и пришли к выводу, что текущие LLM демонстрируют существенный разрыв с нормативной моделью именно в задачах, требующих активного поиска фальсифицирующих свидетельств.

Batista и Griffiths~\cite{batista2026rational} предложили теоретическое объяснение: склонность LLM к подтверждающим тестам может быть рационально обоснована в контексте сикофантических паттернов, усиленных обучением с подкреплением на основе обратной связи от человека. Согласно этой гипотезе, модели, оптимизированные на одобрение человека, могут воспроизводить предпочтение подтверждающей информации как следствие смещения функции вознаграждения.

\subsection{Сравнительные бенчмарки когнитивных предубеждений в LLM}

Ряд работ ставил задачу систематического сравнения нескольких когнитивных задач на одном и том же наборе моделей. Koo и соавторы~\cite{koo2024benchmarking} разработали бенчмарк когнитивных предубеждений для LLM в роли оценщиков, охватывающий несколько типов bias, и обнаружили, что инструктируемые модели воспроизводят человекоподобные предубеждения при оценке качества текстов.

Macmillan-Scott и Musolesi~\cite{macmillan2024irrationality} и Yax и соавторы~\cite{yax2024studying} продемонстрировали, что профиль предубеждений существенно варьируется между моделями и не определяется однозначно размером или семейством. Coda-Forno и соавторы~\cite{coda2024cogbench} создали наиболее полный на момент публикации сравнительный фреймворк, охватывающий задачи из области вероятностного рассуждения, обучения с подкреплением и межвременного выбора.

Настоящее исследование развивает эту линию, фокусируясь на трёх парадигмах, образующих концептуально связный спектр от вероятностного суждения до индуктивного открытия правил, и расширяя охват до девяти моделей с введением оригинальных непрерывных метрик.

% -------------------------------------------------------
\section{Образ решения: исследовательские гипотезы}
% -------------------------------------------------------

\subsection{Гипотеза об универсальности предубеждений}

На основании анализа литературы формулируется первая исследовательская гипотеза.

\textbf{Гипотеза H1 (универсальность).} \textit{Все тестируемые языковые модели демонстрируют статистически значимое предубеждение подтверждения во всех трёх парадигмах относительно нормативного стандарта.}

Основанием для данной гипотезы служат следующие соображения. Обучающие корпуса всех современных LLM преимущественно состоят из текстов, порождённых людьми, в которых человекоподобные паттерны рассуждения, включая систематические ошибки, представлены значительно шире, чем нормативно корректные рассуждения. Кроме того, процедура RLHF (обучения с подкреплением на основе обратной связи от человека) дополнительно усиливает воспроизведение паттернов, одобряемых людьми-оценщиками, которые сами подвержены предубеждениям~\cite{batista2026rational}. Наконец, эмпирические данные предшествующих работ~\cite{binz2023using, macmillan2024irrationality, banatt2024wilt} согласуются с наличием предубеждений у всех протестированных моделей, пусть и с различной интенсивностью.

\subsection{Гипотеза о масштабировании и снижении предубеждений}

\textbf{Гипотеза H2 (масштабирование).} \textit{Более крупные модели демонстрируют более низкий уровень когнитивных предубеждений, однако данная зависимость не является монотонной и может нарушаться внутри модельных семейств.}

Связь между размером модели и когнитивными способностями документирована в ряде работ~\cite{dasgupta2024language, banatt2024wilt}. Однако механизм этой зависимости остаётся дискуссионным: снижение предубеждений при увеличении размера может объясняться как улучшением базовых навыков рассуждения, так и более полным memorization нормативных ответов на классические задачи. Вторая часть гипотезы --- о возможном нарушении монотонности --- мотивирована предварительным анализом данных, выявившим паттерн, при котором меньшие модели одного семейства демонстрируют более низкий уровень предубеждений, чем большие.

\subsection{Гипотеза о спектральной природе предубеждений}

\textbf{Гипотеза H3 (спектральность).} \textit{Интенсивность когнитивных предубеждений систематически возрастает от вероятностной парадигмы к логической и далее к индуктивной, отражая возрастающую сложность требуемого рассуждения.}

Данная гипотеза основана на предположении о том, что более сложные типы рассуждения (многоходовое индуктивное обновление убеждений) в большей мере опираются на эвристические стратегии, чем более простые (одношаговое вероятностное суждение). Следствием является ожидание, что все модели будут демонстрировать наиболее высокий уровень ошибок в задаче 2-4-6 и наиболее низкий~--- в задаче Линды (при контроле за базовым уровнем сложности задач).

\subsection{Гипотеза о демографическом токен-bias}

\textbf{Гипотеза H4 (демографический bias).} \textit{Демографические переменные, кодируемые именем персонажа, оказывают статистически значимое влияние на интенсивность ошибки конъюнкции, при этом различия будут более выражены по оси расы, чем по осям возраста и пола.}

Гипотеза мотивирована методологическими соображениями: если токен-bias действительно определяется ассоциативными связями в обучающем корпусе, то демографические токены, имеющие различную частотность и различные ассоциативные профили, должны порождать различные уровни bias. Дизайн с фиксированными шаблонами и операцией строковой замены~\cite{jiang2024peek} позволяет изолировать эффект демографических токенов от возможных конфаундеров нарративного стиля.

% -------------------------------------------------------
\section{Требования к экспериментальному исследованию}
% -------------------------------------------------------

\subsection{Требования к датасетам и контролю контаминации}

На основании анализа ограничений существующих подходов формулируются следующие требования к датасетам.

\textit{Оригинальность стимулов.} Все биографические виньетки для задачи Линды должны быть составлены вручную и не воспроизводить текст оригинальной задачи или её известных адаптаций. Для задачи Уэйсона и задачи 2-4-6 стимулы должны включать деконтаминированные условия с оригинальным содержанием.

\textit{Структурное единообразие.} Все виньетки и правила одного условия должны быть сконструированы по единому шаблону, исключающему конфаундеры, связанные с различиями в нарративном стиле или синтаксической сложности.

\textit{Размер и балансированность.} Датасеты должны обеспечивать достаточную статистическую мощность для применения используемых тестов. Для теста МакНемара необходимо минимум 25--30 дискордантных пар на каждое сравниваемое условие; для оценки демографических эффектов~--- достаточное число подстановок на каждую расовую группу.

\textit{Верификация.} Каждый элемент датасета должен пройти проверку на отсутствие функциональных дубликатов, тематических пересечений с известными источниками и логических противоречий.

\subsection{Требования к метрикам и статистическому выводу}

\textit{Совместное использование бинарных и непрерывных метрик.} Наряду с долей правильных/ошибочных ответов должна применяться непрерывная метрика, учитывающая уверенность модели в своём выборе. Это обеспечивает обнаружение предубеждений в режиме «потолка точности», когда бинарная метрика не различает модели.

\textit{Парный дизайн и тест МакНемара.} Для сравнения условий (correlated vs. uncorrelated; с перестановкой vs. без) должен применяться тест МакНемара, использующий парную структуру данных и обладающий более высокой статистической мощностью, чем независимые тесты~\cite{mcnemar1947note}.

\textit{Контроль множественных сравнений.} При сравнении нескольких демографических групп или условий промпта должна применяться процедура коррекции на множественные сравнения~\cite{benjamini1995controlling}.

\textit{Воспроизводимость.} Все модели должны запрашиваться с фиксированными параметрами (temperature = 0 для детерминированных условий), а сырые ответы сохраняться отдельно от производных метрик.

\subsection{Требования к набору тестируемых моделей}

Для обеспечения широкой генерализуемости выводов набор тестируемых моделей должен удовлетворять следующим критериям.

\textit{Разнообразие провайдеров и архитектур.} Набор должен включать модели от не менее чем четырёх различных провайдеров, охватывая как проприетарные (OpenAI, Anthropic, xAI), так и открытые (Google, Alibaba, Zhipu) системы.

\textit{Разнообразие размеров.} Набор должен включать модели как уровня «малые» (7--12B параметров), так и «крупные» (70B+ параметров), что позволяет проверить гипотезу о масштабировании.

\textit{Единообразие условий доступа.} Все модели должны запрашиваться через единый API-интерфейс для исключения артефактов, связанных с различием в препроцессинге промптов или постпроцессинге ответов.